\begin{document}
\begin{center}
{\large\bfseries Supplementary material for ``Sensor degradation and maintenance-trigger errors in remaining-useful-life prediction intervals: a C-MAPSS stress test''\par}
\ifanon\else\vspace{0.3em}{\small Yi Wang, Lishui University\par}\fi
\end{center}
\small
\noindent All tables are generated by \texttt{paper\_jrr/make\_tables\_v18.py} (the unchanged v1.8 generator, from released result files) or \texttt{paper\_jrr/make\_tables\_jrr.py} (new or re-derived quantities, from released result files and from \texttt{results\_jrr/derived/}); no value is typed by hand. Rates are percentages of evaluated engines unless stated. Analyses marked \emph{exploratory} were run once after the main protocol was fixed and are not pre-registered.

\section*{S1. Protocol audit}
\textbf{Controlled $2\times2$ retraining.} Table~\ref{tab:leak} trains the heteroscedastic LSTM with the same five seeds under all four combinations of checkpoint selection (test set T or validation engines V) and normaliser scope (whole training file W or fitting engines F), with the fitting engines held fixed. Sel.\ is the mean over normaliser levels of (T$-$V), Norm.\ the mean over selection levels of (W$-$F), Int.\ the difference of differences; negative means the choice makes the result look better. Test-set selection improves RMSE by \NSelEffRMSE{} cycles and the NASA score by \NSelEffScore{} units; the normaliser scope changes RMSE by \NNormEffRMSE{} cycles. The interaction (\NIntEffRMSE{} cycles, \NIntEffScore{} score units) is of the same order as the normaliser effect. The only cell in which a leaked choice does not flatter the result is the \NLeakExceptions. An earlier uncontrolled version left the fitting set free, so whole-file cells were fitted on more engines, and the normaliser appeared to cost as much as the selection choice; that apparent effect was training-set size.
\begin{table}[H]\centering\footnotesize
\caption{Controlled $2\times2$ experiment, heteroscedastic LSTM, terminal unit, mean of five seeds.}\label{tab:leak}
\begin{tabular}{llrrrr|rrr}\toprule
 & & T/W & V/W & V/F & T/F & Sel. & Norm. & Int. \\ \midrule
\tabin{tables/tab_leak.tex}\bottomrule\end{tabular}
\end{table}

\noindent\textbf{Audit history.} The protocol was adopted after an audit of an earlier, unpublished pipeline found that the checkpoint had been selected on the official test set and the normaliser fitted on data used for checkpoint selection. The audit also found an MC Dropout baseline that omitted the residual variance term (coverage 0.40--0.49) and an ensemble evaluation that gave each member its own noise realisation, which averages the perturbation at $1/\sqrt{M}$ and manufactured apparent robustness. All were corrected; the earlier results are retained in the released repository under \texttt{results/superseded/}.

\section*{S2. Clean benchmark, mean processing and ensemble controls}
\noindent\textbf{Predicted-mean processing by analysis.} The analyses of the original pipeline did not post-process the predicted mean identically. Table~\ref{tab:proto} lists what each analysis does, read from the code; the clean baseline of the main text (Table~1) and Table~\ref{tab:samepath} use the stress-test processing.
\begin{table}[H]\centering\scriptsize
\caption{Predicted-mean processing (clip to $[0,125]$ or none) by analysis. Heteroscedastic: mean of the heteroscedastic head, also used by CP-norm. Other: fixed-variance, MC Dropout; ensemble mean.}\label{tab:proto}
\begin{tabular}{p{5.6cm}p{3.6cm}p{4.4cm}}\toprule Analysis (scripts) & Heteroscedastic, CP-norm & Other \\ \midrule
v1.8 clean benchmark: PICP, MPIW, ECE, per-engine, RMSE (\texttt{eval\_*\_b1}, \texttt{eval\_*\_deep\_ensemble}, per-engine coverage) & clipped & clipped; ensemble members clipped before averaging \\
v1.8 clean interval score of the heteroscedastic head, CP-norm, ensemble (\texttt{interval\_score.py}) & unclipped & ensemble members unclipped, ensemble mean clipped \\
v1.8 $2\times2$ retraining, same-split ensemble control (\texttt{protocol\_2x2\_*}, \texttt{*samesplit*}) & clipped & clipped \\
v1.8 sweeps, dose--response, decision tables, drift controls (\texttt{sweep\_engine}, \texttt{maintenance\_decision\_*}, \texttt{drift\_controls}) & unclipped & fixed variance and MC Dropout clipped; ensemble members unclipped, mean clipped \\
Derived analyses of this revision (\texttt{jrr\_decision\_per\_engine}, \texttt{jrr\_coherent\_*}) & unclipped & as the sweeps \\
Clean baseline of this revision (\texttt{jrr\_clean\_same\_path.py}) & unclipped (variant B); clipped (variant A) reported for reference & as the sweeps \\
Clipping sensitivity (\texttt{jrr\_clip\_sensitivity.py}) & both & ensemble members clipped individually (A) or not (B) \\ \bottomrule
\end{tabular}
\end{table}
\begin{table}[H]\centering\scriptsize
\caption{Clean terminal-window benchmark recomputed through the stress-test code path (mean of five seeds; MC Dropout also over five sampling repeats; the ensemble is a single $M=5$ estimate). PICP, MPIW (cycles), ECE, interval score IS and RMSE (cycles) use the unclipped heteroscedastic and CP-norm means; the last column is PICP under the v1.8 benchmark processing (all means clipped).}\label{tab:samepath}
\begin{tabular}{llccccc|c}\toprule Cell & Mechanism & PICP & MPIW & ECE & IS & RMSE & PICP, v1.8 \\ \midrule
\tabin{tables/jrr_supp_clean_same.tex}\bottomrule\end{tabular}
\end{table}
\begin{table}[H]\centering\scriptsize
\caption{Effect of clipping the heteroscedastic mean to $[0,125]$ on coverage (stored per-engine outputs; unclipped / clipped), 90\% interval, terminal window; noise 1\%, bias 5\%, drift 5\% of range. Trigger rates of the heteroscedastic head and CP-norm are unchanged in all \JClipRateN{} cell--condition--lead-time combinations; the ensemble changes by at most \JClipRateMaxEns{} points.}
\begin{tabular}{lcccc}\toprule Cell & Clean & Noise & Bias & Drift \\ \midrule
\tabin{tables/jrr_supp_clip.tex}\bottomrule\end{tabular}
\end{table}
\FloatBarrier
\noindent\textbf{v1.8 clean benchmark (benchmark processing, unchanged).} PICP, MPIW, ECE and per-engine compliance are from the clipped benchmark; the interval score of the heteroscedastic head, CP-norm and ensemble is the unclipped-path value from the interval-score file, and that of the fixed-variance and MC Dropout heads comes from the fair-calibration files, so the IS column mixes processings (compare Table~\ref{tab:samepath}).
\begin{table}[H]\centering\scriptsize
\caption{v1.8 clean-input benchmark under fair calibration, v1.8 processing, terminal unit except per-engine (trajectory unit), mean of five seeds; the ensemble is a single $M=5$ estimate.}
\begin{tabular}{llccccc|ccccc}\toprule
 & & \multicolumn{5}{c|}{LSTM} & \multicolumn{5}{c}{Transformer} \\
 & Mechanism & PICP & MPIW & ECE & Per-eng. & IS & PICP & MPIW & ECE & Per-eng. & IS \\ \midrule
\tabin{tables/tab_clean.tex}\bottomrule\end{tabular}
\end{table}
\begin{table}[H]\centering\scriptsize
\caption{Ensemble gain with the partition held fixed. Single models: five initialisations on the seed-42 partition; same-partition ensemble: those five combined; cross-partition ensemble: the five seeds of the main benchmark, each on its own partition.}
\begin{tabular}{llcccccc}\toprule
 & & \multicolumn{4}{c}{IS} & \multicolumn{2}{c}{PICP} \\
Backbone & Data & 5 part., single & 1 part., single & 1 part., ens. & 5 part., ens. & 1 part., ens. & 5 part., ens. \\ \midrule
\tabin{tables/tab_samesplit.tex}\bottomrule\end{tabular}
\end{table}
The same-partition ensemble beats the mean of its own members by \NSameGain{} IS points in all eight combinations, and the two ensembles differ by at most \NSameCrossMax{} points with no consistent direction. On FD001 five single models on one partition average \NOnePartSingleFDone{} IS against \NFivePartSingleFDone{} for five single models on five partitions, a difference as large as the ensemble margin; five seeds on five partitions are therefore a coarse instrument. Three disjoint five-member ensembles from a fifteen-seed pool give IS of \NIndepIS{} on FD001--FD004; the ratio of fifteen- to five-seed PICP standard deviation is \NNFiveFifteen{} on FD001, FD002 and FD004; over $M\in\{2,3,5,10\}$ the FD004 PICP rises from \NMSweepPICP{} and IS falls from \NMSweepIS.

\noindent\textbf{Fair calibration.} Table~\ref{tab:fair} regenerates the calibration-versus-training residual comparison from released result files. The v1.8 manuscript typed this table by hand; its ratio column agrees with the result files, but its coverage columns did not agree with the main benchmark (differences up to 0.016), and the training-residual coverage is available from the released results only for the LSTM fixed-variance head on FD001, FD002 and FD004.
\begin{table}[H]\centering\footnotesize
\caption{Ratio of calibration-set to training-residual scale, and PICP under each estimator (``--'': not in the released result files).}\label{tab:fair}
\begin{tabular}{lcccc}\toprule
 & Ratio & PICP fixed, train & PICP fixed, calibration & PICP MC Dropout, calibration \\ \midrule
\tabin{tables/jrr_supp_fair.tex}\bottomrule\end{tabular}
\end{table}

\FloatBarrier
\section*{S3. Coverage against out-of-range share}
Perturbation levels are percentages of the training range for noise, bias and drift and percentages of the reading for gain. Matched-excursion interpolation (heteroscedastic head, linear in $f_{\mathrm{oob}}$ between measured levels; ``--'': the target lies outside the measured range of that curve). The underlying points are in \texttt{results\_jrr/derived/foob\_coverage\_points.csv}.
\begin{table}[H]\centering\scriptsize
\caption{Coverage at matched $f_{\mathrm{oob}}$ of 5\% and 10\% (excursion levels, not perturbation amplitudes; amplitudes are percentages of range except gain, which is a percentage of the reading).}
\begin{tabular}{lllcccc}\toprule
Backbone & Data & $f_{\mathrm{oob}}$ & Noise & Bias & Gain & Drift \\ \midrule
\tabin{tables/jrr_supp_matched.tex}\bottomrule\end{tabular}
\end{table}
\begin{table}[H]\centering\scriptsize
\caption{Excursion $f_{\mathrm{oob}}$ at which the heteroscedastic head loses ten points of coverage (v1.8 definition), with left-censored cells marked LC (coverage already lost ten points at the lowest level; the $f_{\mathrm{oob}}$ at that level in parentheses). n.r.: not reached within 5\% of range; $<10^{-4}$: coverage falls while the excursion stays below $10^{-4}$.}
\begin{tabular}{llcccc}\toprule
Backbone & Degradation & FD001 & FD002 & FD003 & FD004 \\ \midrule
\tabin{tables/jrr_supp_types.tex}\bottomrule\end{tabular}
\end{table}

\FloatBarrier
\section*{S4. Decision stress test}
\begin{table}[H]\centering\scriptsize
\caption{Both mechanisms' families at $L=20$ for drift and bias, six-condition sub-datasets; engine-bootstrap 95\% intervals (\%).}
\begin{tabular}{llllcc}\toprule
Cell & Condition & Mechanism & Premature & Unrecognised \\ \midrule
\tabin{tables/jrr_supp_mech.tex}\bottomrule\end{tabular}
\end{table}
\begin{center}\scriptsize
\begin{longtable}{lllcll}
\caption{Heteroscedastic head, all backbones, sub-datasets, conditions and lead times; premature and unrecognised rates (\%) with engine-bootstrap 95\% intervals.}\\
\toprule
Backbone & Data & Condition & $L$ & Premature & Unrecognised \\ \midrule \endfirsthead
\toprule Backbone & Data & Condition & $L$ & Premature & Unrecognised \\ \midrule \endhead
\tabin{tables/jrr_supp_decision_nll.tex}\bottomrule
\end{longtable}
\end{center}
\begin{table}[H]\centering\footnotesize
\caption{Rank stability across cost ratios 5/20/100 by condition: groups (backbone $\times$ sub-dataset $\times$ lead time), groups with a changed full ranking, groups with a changed top place, and the median engine-bootstrap probability that the top mechanism at ratio 5 differs from that at ratio 100 among groups with a changed top place.}
\begin{tabular}{lcccc}\toprule
Condition & Groups & Ranking changes & Top changes & Median probability \\ \midrule
\tabin{tables/jrr_supp_rank.tex}\bottomrule\end{tabular}
\end{table}
\begin{landscape}
\noindent\textbf{Full cost and rank table (v1.8 generator).} Cost $=c_{\mathrm{late}}r_{\mathrm{late}}+r_{\mathrm{early}}$ at ratios 5/20/100 and rank for every mechanism. In this table unrecognised (overall / conditional) and premature rates are fractions between 0 and 1, not percentages; cost is in the same units; RUL@t is the true RUL at premature trigger under the 125-cycle clip. The table is generated by the unchanged v1.8 generator. Rates and engine-bootstrap intervals of all five mechanisms, eight cells, four conditions and three lead times are enumerated in \texttt{results\_jrr/derived/decision\_engine\_bootstrap\_all.csv}.
\begin{center}\scriptsize
\begin{longtable}{lllllcclll}
\toprule Backbone & Data & Cond. & $L$ & Mech. & unrec. (ovr / cond.) & prem. & RUL@t & cost 5/20/100 & rank \\ \midrule \endhead
\tabin{tables/supp_s2.tex}\bottomrule
\end{longtable}
\end{center}\end{landscape}
\noindent Per-engine bootstrap of the top-ranked mechanism against the other four (v1.8; the top is chosen from the same data, so these intervals are descriptive): number of the 144 comparisons per backbone and sub-dataset whose 95\% interval excludes zero.
\begin{center}\footnotesize\begin{tabular}{llc}\toprule Backbone & Data & Excluding zero \\ \midrule
\tabin{tables/supp_s2c.tex}\bottomrule\end{tabular}\end{center}

\FloatBarrier
\section*{S5. Output swap, controls and additive-noise attribution}
\begin{landscape}
\begin{table}[H]\centering\scriptsize
\caption{Output swap at the decision point, heteroscedastic head, $L=20$ (exploratory). $\Delta\mu$: mean shift in cycles (perturbed minus clean, same model). Premature rates: clean, mean perturbed only, scale perturbed only, both; $\Delta_\sigma$: effect of the perturbed scale given the perturbed mean (points, with engine-bootstrap interval). Unrecognised rates: mean only, scale only, both.}
\begin{tabular}{llcccccccccc}\toprule
Cell & Cond. & $\Delta\mu$ & Prem. clean & $\mu$ only & $\sigma$ only & both & $\Delta_\sigma\mid\mu_1$ & Unrec. $\mu$ only & $\sigma$ only & both \\ \midrule
\tabin{tables/jrr_supp_swap.tex}\bottomrule\end{tabular}
\end{table}
\begin{table}[H]\centering\scriptsize
\caption{Output swap for the sign-pattern variants of the 5\% shift (Section~3.3), heteroscedastic head, $L=20$, exploratory, inference only. $\Delta\mu$: mean shift in cycles; $\sigma$/$\sigma_0$: mean ratio of perturbed to clean scale. Effects are in percentage points with engine-bootstrap 95\% intervals: mean effect = perturbed mean with clean scale minus clean; scale effect = both perturbed minus perturbed mean with clean scale.}\label{tab:cohswap}
\begin{tabular}{llccccccc}\toprule
Cell & Shift & $\Delta\mu$ & $\sigma/\sigma_0$ & Prem., mean effect & Prem., scale effect & Unrec., mean effect & Unrec., scale effect \\ \midrule
\tabin{tables/jrr_supp_cohswap.tex}\bottomrule\end{tabular}
\end{table}
\end{landscape}
\begin{table}[H]\centering\scriptsize
\caption{Waveform controls at 5\% of range, heteroscedastic head, terminal-window PICP (released result files except the constant $+$ offset, which is from the sign-coherent control). Ramp: drift family; reversed, shuffled, per-trajectory (continuous), bias with equal end value and single-channel ramp as in the earlier analysis; $k=1,3,5$: fixed-end shuffled ramp (FD002 only; ``--'': not available).}\label{tab:wf}
\begin{tabular}{lcccccccc|ccc}\toprule
Cell & Clean & Ramp & Const.\ $+$ & Reversed & Shuffled & Per traj. & Bias eq. & Single ch. & $k=1$ & $k=3$ & $k=5$ \\ \midrule
\tabin{tables/jrr_supp_waveform.tex}\bottomrule\end{tabular}
\end{table}
\begin{table}[H]\centering\footnotesize
\caption{Drift controls at 5\% of range, FD002, heteroscedastic head: excursion and PICP.}
\begin{tabular}{lccc}\toprule Perturbation & $f_{\mathrm{oob}}$ & LSTM & Transformer \\ \midrule
\tabin{tables/tab_controls.tex}\bottomrule\end{tabular}
\end{table}
\noindent\textbf{Additive-noise attribution (v1.8 tables).} Four-way frozen-output decomposition at the point where the heteroscedastic head has lost ten points of coverage. Scale contribution: share of the total ten-point loss (positive: the live scale makes the loss larger), under both decomposition orders.
\begin{table}[H]\centering\footnotesize
\caption{Attribution, LSTM.}\begin{tabular}{llcccc}\toprule & & FD001 & FD002 & FD003 & FD004 \\ \midrule
\tabin{tables/tab_attr.tex}\bottomrule\end{tabular}
\end{table}
\begin{table}[H]\centering\footnotesize
\caption{Attribution, Transformer.}\begin{tabular}{llcccc}\toprule & & FD001 & FD002 & FD003 & FD004 \\ \midrule
\tabin{tables/tab_attr_transformer.tex}\bottomrule\end{tabular}
\end{table}
\begin{landscape}\begin{center}\scriptsize
\begin{longtable}{lllccccccccccc}
\caption{Four-way attribution values; $C_{ij}$: coverage with mean from condition $i$ and scale from condition $j$ (0 clean, 1 perturbed); order A varies the mean first.}\\
\toprule Backbone & Data & Mech. & $f_{\mathrm{oob}}$ (\%) & $C_{00}$ & $C_{10}$ & $C_{01}$ & $C_{11}$ & $\mu$ (A) & $\sigma$ (A) & $\mu$ (B) & $\sigma$ (B) & Int. & Both \\ \midrule \endfirsthead
\toprule Backbone & Data & Mech. & $f_{\mathrm{oob}}$ (\%) & $C_{00}$ & $C_{10}$ & $C_{01}$ & $C_{11}$ & $\mu$ (A) & $\sigma$ (A) & $\mu$ (B) & $\sigma$ (B) & Int. & Both \\ \midrule \endhead
\tabin{tables/supp_s1.tex}\bottomrule
\end{longtable}
\end{center}\end{landscape}
\begin{center}\scriptsize\begin{longtable}{lllcccc}
\caption{Per-seed paired bootstrap of $|\Delta_\mu|-|\Delta_\sigma|$ (five seeds, 2000 resamples; conditional on the test set, the seeded realisations and the interpolated crossover).}\\
\toprule Backbone & Data & Mech. & Mean first: mean [95\% CI] & $>0$ & Scale first: mean [95\% CI] & $>0$ \\ \midrule \endhead
\tabin{tables/supp_boot.tex}\bottomrule
\end{longtable}\end{center}

\FloatBarrier
\section*{S6. Sensitivity and scope checks}
\noindent\textbf{Evaluated engines.} FD001--FD004: \JEvalEng{} evaluated against \JFileEng{} in the test files; engines with fewer than 30 recorded cycles cannot form a window. \textbf{Label instant.} Scoring against the current-cycle label (official RUL, clipped) instead of $\min(\mathrm{RUL}-1,125)$ changes any decision rate of any mechanism, condition, lead time or backbone by at most \JLabelMax{} percentage points; for the headline drift and bias cells it changes them by at most \JLabelHeadline{} points. \textbf{Operating settings.} Perturbing the three operating-setting channels as well as the sensors raised the premature rate under drift on FD004 from \NAOneFDfourSens{} to \NAOneFDfourJoint; the settings alone gave \NAOneFDfourSet{}; on FD002 \NAOneFDtwoSens{} against \NAOneFDtwoJoint. \textbf{Agreement with the release.} The per-engine re-derivation reproduced \JAgreeN{} released decision rates with maximum absolute difference \JAgreeMax{}, and the released ramp coverage for all eight backbone--sub-dataset cells (\texttt{coherent\_bias\_control.json}).
\begin{table}[H]\centering\scriptsize
\caption{Sign-coherent shift control, heteroscedastic head, 5\% of range: PICP / premature (\%) / unrecognised (\%), $L=20$ (exploratory).}
\begin{tabular}{lcccc}\toprule & Clean PICP & Ramp (+) & Constant $+5\%$ & Constant $-5\%$ \\ \midrule
\tabin{tables/jrr_tab_coh.tex}\bottomrule\end{tabular}
\end{table}

\FloatBarrier
\section*{S7. Provenance, latency and reference verification}
\begin{center}\scriptsize
\begin{longtable}{p{2.0cm}p{3.8cm}p{1.4cm}p{1.4cm}p{0.9cm}p{2.2cm}p{1.2cm}}
\caption{Provenance of every v1.8 table: script or result file, label version, evaluation unit, engine weighting, aggregation, perturbation-scale reference.}\\
\toprule Table & Script / result file & Labels & Unit & Eng.-wt. & Aggregation & Noise range \\ \midrule \endhead
\tabin{tables/supp_s0.tex}\bottomrule
\end{longtable}
\end{center}
{\sloppy \noindent New analyses of this revision: \path{code_jrr/jrr_decision_per_engine.py} (inference only; per-engine lower bounds), \path{code_jrr/jrr_derived_stats.py} (engine bootstrap, output swap, label-instant sensitivity, rank stability, $f_{\mathrm{oob}}$ points), \path{code_jrr/jrr_coherent_bias_control.py} (sign-coherent shift control), \path{code_jrr/jrr_coherent_swap.py} (output swap for the sign-pattern variants), \path{code_jrr/jrr_clean_same_path.py} (clean baseline under the stress-test processing) and \path{code_jrr/jrr_clip_sensitivity.py} (clipping sensitivity). The CSV files \path{decision_engine_bootstrap_all.csv}, \path{rank_stability_all96.csv} and \path{clean_same_path_table.csv} in \path{results_jrr/derived/} enumerate the corresponding tables in full. Outputs are under \path{results_jrr/derived/}; \path{results/} is unchanged from the original release\ifanon\else{} (tag \path{v1.8-submission})\fi.}

\noindent\textbf{Latency (v1.8 measurement, not re-run).} Per-sample latency in microseconds, median (IQR) over eleven repetitions, one NVIDIA RTX A5000 Laptop GPU, CUDA-event timing; columns: heteroscedastic, fixed variance, CP-norm, MC Dropout ($T=50$), ensemble ($M=5$).
\begin{center}\scriptsize\begin{longtable}{lllccccc}\toprule Data & Backbone & Batch & Heterosc. & Fixed var. & CP-norm & MC Drop. & Ensemble \\ \midrule \endhead
\tabin{tables/supp_s4_cuda.tex}\bottomrule\end{longtable}\end{center}
\noindent\textbf{Reference verification.} Each reference listed in the manuscript was checked against Crossref (DOI and title) or, for the preprint, the arXiv abstract page, on 12 October 2026. References without DOI in the source bibliography were not machine-checked.
\begin{center}\scriptsize\begin{longtable}{llp{9cm}}\toprule \# & Key & Check \\ \midrule \endhead
\tabin{tables/jrr_supp_refcheck.tex}\bottomrule\end{longtable}\end{center}
\end{document}
